\section{Pascal: the universal identity and its constructive consequence}\label{sec:pascal}
Write $[ijk]=\det(P_i,P_j,P_k)$, with vectors as columns, and put
\[
 X=(P_1\times P_2)\times(P_4\times P_5),\quad
 Y=(P_2\times P_3)\times(P_5\times P_6),\quad
 Z=(P_3\times P_4)\times(P_6\times P_1).
\]
For $P=(x,y,z)$ let
$\nu(P)=(x^2,y^2,z^2,xy,xz,yz)$.  The universal identity is
\begin{equation}\label{eq:pascal-universal}
 \det(X,Y,Z)=\det\begin{pmatrix}
 \nu(P_1)\\\nu(P_2)\\\nu(P_3)\\
 \nu(P_4)\\\nu(P_5)\\\nu(P_6)
 \end{pmatrix}.
\end{equation}
This is an identity over $\mathbb Z$ in eighteen indeterminates.  Consequently
it remains valid after substitution into every commutative ring, including
rings of characteristic two.  The following proof never assumes that a
particular triple of input points is independent.

\subsection{The determinant relations used in the proof}
For five vectors $r,s,t,u,v$, the bracket relation is
\begin{equation}\label{eq:basic-plucker}
 [rst][ruv]-[rsu][rtv]+[rsv][rtu]=0.
\end{equation}
Here the letter bracket denotes the same three-column determinant.
To prove it without a nondegeneracy hypothesis, first expand the determinant
of the four-by-four matrix whose last three rows are the coordinates of
$r,s,t,u$ and whose first row repeats any one of those three rows.  Since
the matrix has a repeated row, its determinant is zero.  Expansion in the
first row, for each of its three possible choices, gives the vector identity
\[
 [stu]r-[rtu]s+[rsu]t-[rst]u=0.
\]
Taking its scalar product with $r\times v$ proves
\eqref{eq:basic-plucker}.  This proves the relation over every commutative
ring by finite determinant algebra.

The four particular instances needed below, with their signs displayed, are
\begin{align}
 [125][146]-[124][156]&=[126][145],\\
 [126][235]-[125][236]&=-[123][256],\\
 [134][236]-[234][136]&=[123][346],\\
 [245][156]-[145][256]&=-[125][456].
\end{align}
They result from taking respectively the common first vector in
\eqref{eq:basic-plucker} to be $P_1,P_2,P_3,P_5$ and reordering the other
arguments.  No division is involved.

\subsection{A two-term bracket expression}
The double-cross-product identity gives
\[
 X=[125]P_4-[124]P_5,\quad
 Y=[236]P_5-[235]P_6,\quad
 Z=[134]P_6-[346]P_1.
\]
On expanding the determinant, four of its eight terms vanish because they
have repeated columns.  Thus, writing $F=\det(X,Y,Z)$,
\begin{align*}
 F={}&[125][236][134][456]
       -[125][236][346][145]\\
    &+[125][235][346][146]
       -[124][235][346][156].
\end{align*}
The last two terms combine, by the first displayed bracket relation, to
$[235][346][126][145]$.  Combining that result with the second term and
using the second bracket relation yields
\begin{equation}\label{eq:pascal-two-term}
 F=[125][134][236][456]-[123][145][256][346].
\end{equation}

\subsection{Why the polynomial is alternating}
Interchanging $P_1$ and $P_2$ in \eqref{eq:pascal-two-term} and adding the
result to the original polynomial gives
\begin{align*}
 F+F_{(12)}={}&[125][456]\bigl([134][236]-[234][136]\bigr)\\
 &+[123][346]\bigl([245][156]-[145][256]\bigr)=0,
\end{align*}
by the third and fourth bracket relations.  Therefore $F$ changes sign
under the transposition $(12)$.

A cyclic shift of the six labels sends $(X,Y,Z)$ to $(Y,Z,-X)$.
Its effect on $F$ is consequently multiplication by $-1$.
Conjugating $(12)$ by powers of this cycle gives every adjacent
transposition; the two signs contributed by the conjugating powers cancel.
Thus $F$ changes sign under each adjacent transposition, and hence under any
permutation by that permutation's sign.  This is an exact polynomial
argument, with no appeal to a dense set or to a limiting configuration.

Each vertex occurs twice in every term of $F$.  Hence $F$ is homogeneous of
degree two in the three coordinates of each $P_i$ separately.  Each such
block has exactly the six possible monomials listed in $\nu(P_i)$.
Expand $F$ in products
\[
 \prod_{i=1}^6 \nu(P_i)_{k_i},\qquad k_i\in\{1,\ldots,6\}.
\]
If two indices $k_i,k_j$ are equal, swapping those two vertex blocks leaves
this monomial unchanged, whereas it negates its coefficient.  An integer
coefficient equal to its own negative is zero.  Thus every surviving tuple
$(k_1,\ldots,k_6)$ is a permutation of $(1,\ldots,6)$.
Antisymmetry further shows that its coefficient is the coefficient for
$(1,\ldots,6)$ multiplied by the sign of that permutation.  This is exactly
the determinant expansion.  Therefore $F=c\det(\nu(P_i))$ for one integer
$c$.

To determine $c$, specialize
\[
 P_1=(1,0,0),\quad P_2=(0,1,0),\quad P_3=(0,0,1),
\]
\[
 P_4=(u,v,w),\quad P_5=(x,y,z),\quad P_6=(p,q,r).
\]
Then
\[
 X=(uz-wx,vz-wy,0),\quad
 Y=(0,py-qx,pz-rx),\quad
 Z=(-qu,-qv,-rv).
\]
Their determinant is the following six-term polynomial:
\begin{align*}
 F={}&pquwyz-pqvwxz-pruvyz\\
     &+prvwxy+qruvxz-qruwxy.
\end{align*}
The coefficient of $(uv)(xz)(qr)$ is $+1$.
In the Veronese determinant the first three rows are the first three basis
rows, and its remaining minor has rows
$(uv,uw,vw)$, $(xy,xz,yz)$, $(pq,pr,qr)$; the same monomial is its diagonal
term, also of coefficient $+1$.  Hence $c=1$, proving
\eqref{eq:pascal-universal}.

The temporary cancellation of $2c=0$ in this proof occurs only in the
integer coefficient ring.  Once the resulting polynomial identity is
proved over $\mathbb Z$, specialization proves it in arbitrary commutative
rings.  In particular, this argument has not imposed a characteristic
restriction on the identity itself.

\subsection{From a conic equation to Pascal collinearity}
Let the six integer vectors satisfy
\[
 Q(P_i)=0\quad(1\leq i\leq6),\qquad
 Q(x,y,z)=q_1x^2+q_2y^2+q_3z^2+q_4xy+q_5xz+q_6yz,
\]
and suppose one specified coefficient $q_j$ is nonzero.  Put
$M_{ik}=\nu(P_i)_k$ and let $C_{ij}$ be its $(i,j)$ cofactor.
For each $j$ the completely explicit identity
\begin{equation}\label{eq:conic-certificate}
 q_j\det M=\sum_{i=1}^6 C_{ij}Q(P_i)
\end{equation}
holds.  Indeed, expand the right side in the coefficients $q_k$.
The coefficient of $q_k$ is $\sum_i C_{ij}M_{ik}$: this is $\det M$ when
$k=j$, by expansion down column $j$, and zero otherwise, by expansion of
the determinant obtained by replacing column $j$ with column $k$, which
has two identical columns.  The hypotheses make the right side zero.
Since $q_j\ne0$ and integer multiplication has no zero divisors,
$\det M=0$.  Equation~\eqref{eq:pascal-universal} gives
$\det(X,Y,Z)=0$.

This proves the finite-coordinate Pascal theorem even for a nonzero
singular conic.  Nonsingularity is not needed for the determinant
conclusion.  To read $X,Y,Z$ as actual projective intersection points,
the relevant adjacent vertices must define nonzero side lines and the
opposite side pairs must be distinct, so that the three double cross
products are nonzero.  The determinant statement remains valid when
these geometric objects degenerate, but the zero vector is not a
projective point. When $X,Y,Z$ are nonzero, Lemma~\ref{lem:axis}
constructs the actual nonzero common line $\axis(X,Y,Z)$. This completes
the proof of Theorem~\ref{thm:pas-summary}, including its witness.

Every rational configuration has integer homogeneous representatives:
multiply each point and the conic coefficients by the product of its
finitely many positive denominators.  This is a primitive-recursive
operation; no gcd, search for a root, or limiting process is required.
The same proof applies over a commutative integral domain, and over a
constructive field when the selected coefficient has an apartness witness.
For the literal first-order iPRA theorem, however, the integer-coded
formulation above is the one being asserted.

\subsection{An exact example with a point at infinity}
For the conic $xz-y^2=0$, take $P(t)=(t^2,t,1)$ for
$t=0,1,2,3,4,5$ in that order.  A chord through parameters $s,t$ has
representative $(1,-(s+t),st)$.  The three opposite-side intersections have
representatives
\[
 (2,2,1),\qquad(7,3,1),\qquad(5,1,0).
\]
They all satisfy the line equation $-x+5y-8z=0$.
The last point lies at infinity, which illustrates why homogeneous
coordinates remove any need for an exceptional ``parallel lines'' case.

\subsection{Verification artifact}
The dependency-free program \texttt{verify\_pascal.py} implements finite
integer-coefficient polynomials as dictionaries of sorted monomials.
It checks the universal determinant identity, the two-term bracket
formula, all four displayed Pluecker instances, and all six cofactor
certificates \eqref{eq:conic-certificate} with indeterminate conic
coefficients.  The Pascal polynomial has exactly $720$ nonzero monomials,
each with coefficient $+1$ or $-1$.  The program performs exact polynomial
normalization, not random sampling.  This is exact computational evidence; it is not a compiled Lean proof or
an emitted iPRA proof object.
